\documentclass{report}
\usepackage{amsmath,amssymb}
\setlength{\parindent}{0mm}
\setlength{\parskip}{1em}
\begin{document}
\begin{center}
\rule{6in}{1pt} \
{\large Deena Hannoun \\
{\bf Simulating Non-Dilute Transport in Porous Media Using a TCAT-Based Model}}

3212 Enchanting Way \\ Raleigh \\ NC 27616
\\
{\tt dhannou@ncsu.edu}\\
Pamela Birak\\
William Gray\\
Tim Kelley\\
Casey Miller\end{center}

Predicting the transport of non-dilute species in fluids of variable
density in porous media is a challenging problem for which existing
mathematical models are unable to represent accurately the experimental
data collected to date. In this work, we consider the displacement of an
aqueous phase with dense brine solutions containing a concentrated CaBr2
species, which in the most concentrated case had a density of 1.71
g/cm$^3$ and a dynamic viscosity of 0.058 cp. Displacement experiments
were conducted in vertically oriented one-dimensional columns for stable
displacements of one fluid by a more dense fluid. Simulation of a
non-dilute system based upon the thermodynamically constrained averaging
theory (TCAT) using an entity-based momentum equation was compared to the
data collected. The model accounts for the effects of non-dilute,
non-ideal systems and consists of a nonlinear set of equations including
a flow equation, a species transport equation, and closure relations. We
rewrite the TCAT entity-based model as a system of two coupled partial
differential-algebraic equations with the relevant closure relations. We
then use a stiff temporal integrator to create 1D simulations of the
model. The model is nonlinear and nonsmooth. We will discuss both results
and numerical difficulties.


\end{document}
